\section{Train Time y Test Time: qué componentes se entrenan realmente}

La forma más eficaz de entender los artículos sobre RAG no es memorizar los nombres de los modelos por año, sino plantear dos preguntas: qué se actualiza durante el entrenamiento y qué se ejecuta durante la prueba. Las diferencias centrales entre Frozen RAG, RePlug, RAG, REALM y Atlas pueden ubicarse en los límites de actualización de language model, query encoder, document encoder, reranker, index y prompt.

\subsection{Entrenamiento e inferencia se registran por separado}

En el lado izquierdo, la slide Train Time and Test Time pregunta si se actualiza el LM, el query encoder o el document encoder, si se actualiza todo a la vez y si se preentrena desde cero; en el lado derecho pregunta si el sistema es un frozen model, cómo cambia la retrieval y si distintos ejemplos usan un index distinto o el mismo.

\lecturefigure{slide-10-train-test-time.jpg}{Matriz de decisiones de RAG en el entrenamiento y en la inferencia}{Video de la clase de Stanford 00:10:54}

\begin{knowledgebox}{Lectura de la figura: un mismo componente puede estar congelado en el entrenamiento y seguir participando en la inferencia}
``Frozen'' significa que los parámetros no se actualizan, no que el componente deje de ejecutarse. Un frozen retriever sigue buscando y un frozen generator sigue generando; el index también puede actualizar sus documentos sin que se entrene el encoder. Al comparar artículos hay que indicar a la vez parameter update, index refresh e inference calls.
\end{knowledgebox}

\begin{importantbox}{Separación en capas entre parámetros y estado}
Sea el conjunto de parámetros
\[
\Theta=\{\theta_G,\theta_Q,\theta_D,\theta_R\},
\]
que representan, respectivamente, el generator, el query encoder, el document encoder y el reranker; sea $I_t$ el estado del index en el instante $t$. El entrenamiento puede actualizar un subconjunto de $\Theta$, y el despliegue puede limitarse a sustituir $I_t$. El aprendizaje de parámetros y la actualización de la base de conocimiento son dos ciclos de vida distintos.
\end{importantbox}

\subsection{Frozen RAG: una línea base sólida sin entrenamiento}

La slide de Frozen RAG agrupa no training e in-context learning: se usa un retriever o una vector database ya existentes para encontrar documentos, se insertan los documentos en el prompt y después se llama a un LLM congelado. Se implementa rápido y es adecuado para black-box API, pero el retriever y el generator no están optimizados el uno para el otro.

\lecturefigure{slide-11-frozen-rag.jpg}{Frozen RAG: el retriever y el LLM se acoplan de forma laxa mediante el prompt}{Video de la clase de Stanford 00:11:57}

\begin{knowledgebox}{Lectura de la figura: valor de ingeniería y límites de investigación de Frozen RAG}
Su valor está en que los módulos son reemplazables, no requiere los pesos del modelo y permite conectar documentos privados con rapidez; sus límites son que el retrieval score no necesariamente coincide con la generator utility, que el orden de los chunks y su posición en el prompt pueden no corresponder, y que el generator puede apoyarse en su memoria paramétrica por encima de la evidencia.
\end{knowledgebox}

\begin{warningbox}{Frozen no significa barato}
Cada solicitud sigue requiriendo embedding, search, reranking, prompt tokens y generation. Si top-$k$ es demasiado grande, el costo de los token de entrada y la latency pueden superar el beneficio de la recuperación; si top-$k$ es demasiado pequeño, puede perderse evidencia clave.
\end{warningbox}

\teachervoice{Kiela llama a frozen RAG ``In-Context Learning only'' y señala repetidamente que es un punto de partida práctico muy sólido, pero también un Frankenstein system cuyos componentes se entrenaron por separado.}

\subsection{Contextualization via Retrieval: qué lado entrenar primero}

La siguiente architecture slide resalta con un recuadro rojo los documents, el document encoder, el retriever y el query encoder, lo que indica que primero se adapta el retrieval side a la tarea. El generator puede seguir congelado, y la señal de entrenamiento proviene de la supervisión de la recuperación, de la likelihood de la respuesta final o de un reranking objective.

\lecturefigure{slide-12-contextualization-via-retrieval.jpg}{Ruta arquitectónica que primero contextualiza el retrieval side}{Video de la clase de Stanford 00:12:09}

\begin{knowledgebox}{Tres tipos comunes de retriever supervision}
\begin{enumerate}
  \item \textbf{Direct relevance labels}: query--positive passage pairs y hard negatives;
  \item \textbf{Answer-containing heuristic}: si el documento contiene la cadena de la respuesta; es barata pero tiene mucho ruido;
  \item \textbf{Generator-derived signal}: qué documento aumenta la likelihood de la respuesta correcta; se acerca más a la tarea final, pero su cómputo es más costoso.
\end{enumerate}
\end{knowledgebox}

\teachervoice{La taxonomy de la clase no es una elección binaria entre «entrenar o no entrenar», sino una revisión elemento por elemento de LM, query encoder, document encoder, reranker e index. Todos los modelos posteriores pueden ubicarse de nuevo en esta tabla.}

\subsection{Resumen de la sección}

Una arquitectura RAG debe reportar a la vez parameter updates, index refresh e inference behavior. Frozen RAG obtiene capacidad rápidamente con un entrenamiento mínimo; los métodos con acoplamiento más estrecho hacen que el retriever o el generator reciban la señal de la tarea final, a costa de mayor complejidad de entrenamiento, actualización del índice y costo de cómputo.
